% from ega0-21-fr.tex
\oldpage[0\textsubscript{IV-1}]{167}
\subsection{Corps admissibles pour une extension}
\label{section:0.21.6-fr}

\begin{env}[21.6.1]
\label{0.21.6.1-fr}
Étant donnés quatre corps $k_0\subset k\subset K\subset L$, il résulte de
(20.6.16) et (20.6.17) que l'on a une suite exacte
\begin{equation}
\tag{21.6.1.1}
\label{0.21.6.1.1-fr}
0\to\Upsilon_{K/k/k_0}\otimes_KL\xrightarrow{v}\Upsilon_{L/k/k_0}
\xrightarrow{u}\Upsilon_{L/K/k_0}\xrightarrow{s}\Upsilon_{L/K/k}\to0.
\end{equation}

Lorsqu'on laisse fixes $k_0$, $K$ et $L$ et qu'on fait «varier» le corps
intermédiaire $k$ entre $k_0$ et $K$, on a évidemment
$\Upsilon_{L/K/k_0}=\Upsilon_{L/K/k}$ lorsque $k=k_0$. Lorsque
l'homomorphisme canonique $s$ de (21.6.1.1) est encore bijectif, on dit que
$k$ est un corps $k_0$-admissible pour l'extension $L$ de $K$. L'intérêt de
l'existence, sous certaines conditions, de tels corps $k$, qui soient pourtant
«suffisamment proches» de $K$ (par exemple tels que $[K:k]$ soit fini), est
qu'ils permettent de remplacer dans certaines questions les modules de
différentielles $\Omega^1_{K/k_0}$ et $\Omega^1_{L/k_0}$ (qui peuvent être
«trop grands», par exemple lorsque $k_0$ est le corps premier) par
$\Omega^1_{K/k}$ et $\Omega^1_{L/k}$, plus aisément maniables.

Lorsque $k_0$ est le \emph{corps premier}, on dira «corps admissible» au lieu
de «corps $k_0$-admissible».
 
On posera
\end{env}

\begin{env}[21.6.1.2]
\label{0.21.6.1.2-fr}
\[
\Delta(L/K,k/k_0)=\operatorname{Coker}
(\Upsilon_{K/k/k_0}\otimes_KL\to\Upsilon_{L/k/k_0})
\simeq\operatorname{Ker}(\Upsilon_{L/K/k_0}\to\Upsilon_{L/K/k})
\]
(espace vectoriel sur $L$)~; son rang sera noté $d(L/K,k/k_0)$ et appelé
\emph{défaut de $k_0$-admissibilité de $k$ pour l'extension $L$ de $K$}
(il est évidemment nul si et seulement si $k$ est $k_0$-admissible pour cette
extension). Lorsque $k_0$ est le corps premier, on écrira
$\Delta(L/K,k)$ et $d(L/K,k)$ au lieu de $\Delta(L/K,k/k_0)$ et
$d(L/K,k/k_0)$.
\end{env}

\begin{proposition}[21.6.2]
\label{0.21.6.2-fr}
Soient $k_0\subset k\subset K\subset L$ quatre corps.
\begin{enumerate}
  \item[\rm (i)] Les conditions suivantes sont équivalentes~:
  \begin{enumerate}
    \item[\rm a)] Le corps $k$ est $k_0$-admissible pour l'extension $L$ de $K$
    (autrement dit, l'homomorphisme $\Upsilon_{L/K/k_0}\to
    \Upsilon_{L/K/k}$ est injectif, donc bijectif).
    \item[\rm b)] L'homomorphisme canonique
    $u:\Upsilon_{K/k/k_0}\to\Upsilon_{L/k/k_0}$ est nul.
    \item[\rm c)] L'homomorphisme canonique
    $v:\Upsilon_{K/k/k_0}\otimes_KL\to\Upsilon_{L/k/k_0}$ est surjectif
    (donc bijectif).
    \item[\rm d)] On a $d(L/K,k/k_0)=0$ (ou
    $\Delta(L/K,k/k_0)=0$).
  \end{enumerate}
  \item[\rm (ii)] Les conditions équivalentes de (i) sont vérifiées lorsque
  l'on est dans l'un des cas suivants~: $\alpha)$ $L$ est séparable sur $k$~;
  $\beta)$ $L$ est séparable sur $K$~; $\gamma)$ on a
  $k\subset k_0(K^p)$, en désignant par $p$ l'exposant caractéristique de $k_0$.
\end{enumerate}
\end{proposition}

\begin{proof}
\begin{enumerate}
  \item[(i)] Les assertions résultent trivialement de l'exactitude de la suite
  (21.6.1.1).
  \item[(ii)] Si $L$ est séparable sur $k$, on a
  $\Upsilon_{L/K/k_0}=0$ (20.6.19), donc la condition b) de (i) est remplie~;
  si $L$ est séparable sur $K$, on a
  $\Upsilon_{L/K/k_0}=0$ (20.6.19), donc la condition a) de (i) est remplie~;
  enfin, si l'on a $k\subset k_0(K^p)$, il en résulte que
  $\Upsilon_{L/K/k_0}=\Upsilon_{L/K/k}$ en vertu de (21.1.5.1), et la
  condition a) de (i) est vérifiée.
\end{enumerate}
\end{proof}
